% Scope: Build FID, spin echo, gradient echo, spoiled GRE and inversion recovery from explicit magnetization histories; establish the sequence vocabulary used later.
% Status: architecture only; no chapter prose or completed problems yet.
% Prerequisites: Chapters 3 through 6.
% Planned worked examples: Spin-echo signal, spoiled-GRE signal, Ernst angle and inversion null time.
% Planned figures: FID, Hahn echo, GRE and inversion-recovery timing diagrams.
% Planned challenge problems: Infer a sequence from timing; optimize flip angle with competing tissue contrasts.
\chapter{Basic Pulse Sequences}
\label{ch:07}

\section{Reading a sequence diagram}
\label{sec:07:reading-a-sequence-diagram}
\subsection{RF, gradient and ADC event tracks}
\subsection{TR, TE, TI and timing references}
\subsection{Two-dimensional and three-dimensional acquisition loops}

\section{Free induction and Hahn spin echo}
\label{sec:07:free-induction-and-hahn-spin-echo}
\subsection{FID and static phase dispersion}
\subsection{Refocusing and echo formation}
\subsection{What a spin echo does and does not reverse}

\section{Gradient echoes}
\label{sec:07:gradient-echoes}
\subsection{Gradient reversal and echo time}
\subsection{T2-star weighting and signal phase}
\subsection{Comparison with RF refocusing}

\section{Spoiled gradient echo}
\label{sec:07:spoiled-gradient-echo}
\subsection{Longitudinal recurrence and steady state}
\subsection{FLASH and SPGR families}
\subsection{Ernst angle and the limits of perfect spoiling}

\section{Inversion and saturation recovery}
\label{sec:07:inversion-and-saturation-recovery}
\subsection{Finite-TR signal equations}
\subsection{Null times and imperfect inversion}
\subsection{STIR, FLAIR and contrast-selective suppression}

\section{Sequence comparisons}
\label{sec:07:sequence-comparisons}
\subsection{Transient versus steady-state acquisition}
\subsection{SNR efficiency and scan-time accounting}
\subsection{Clinical examples and hidden model assumptions}

\section{Worked examples}
\label{sec:07:worked-examples}
% Planned calculations: Spin-echo signal, spoiled-GRE signal, Ernst angle and inversion null time.
\subsection{Derivation and quantitative calculation}
\subsection{Assumptions, limiting cases and scanner interpretation}

\section{Challenge problems}
\label{sec:07:challenge-problems}
% Problem design: Infer a sequence from timing; optimize flip angle with competing tissue contrasts.
% Use challengeproblem and stable labels prob:07:<topic>.
% Complete solutions belong in Appendix E, section for this chapter.
% Use \begin{solution}{prob:07:<topic>} ... \end{solution} there.
\subsection{Derivation and quantitative design}
\subsection{Model criticism and integrated reasoning}
